\section{Secondary Diagnostics and Heterogeneity}

Boundary behavior was substantial in both models. At least one parameter was at a bound for 422 of 617 \QL{} fits (68.40\%) and 393 of 617 \PVL{} fits (63.70\%). At the parameter level, notable accumulation occurred at both learning-rate endpoints in \QL{}, at the upper inverse-temperature bound, and at several \PVL{} bounds. These findings motivate caution when interpreting individual parameter estimates as direct psychological traits.

\begin{figure}[H]
    \centering
    \includegraphics[width=0.82\textwidth]{boundary_fit_rates.pdf}
    \caption{Proportion of fitted participants with parameter estimates at optimization bounds.}\label{fig:boundary-rates}
\end{figure}

Participants came from ten source studies. Model preference was therefore also summarized separately by study. These summaries provide a heterogeneity check, but the inferential analysis treats participants as independent observational units and does not fit a hierarchical study-level model.

\begin{figure}[H]
    \centering
    \includegraphics[width=0.9\textwidth]{study_pvl_preference_rates.pdf}
    \caption{PVL-Delta preference rates by source study for AIC and BIC.}\label{fig:study-preference}
\end{figure}
